% /root/monodromy, 2026-11-19. Complete independent parent proof audit passed.
\subsection{Global source smoothness in the symbolic contact slice}
\label{subsec:ss-uniform-source-smoothness}
\status{Proved and independently audited. The parent checked all chart
incidences, the four exact discriminant multiplicities and the degree
argument. This addresses source smoothness. Other collisions for odd graph
kernels are excluded in Proposition~\ref{prop:ss-uniform-graph-mixed-contact}.
The later uniform tangent-triple and odd-degree full-image results are
Proposition~\ref{prop:ss-uniform-tangent-triple} and
Theorem~\ref{thm:ss-all-odd-full-groups}.}

Keep the notation of Lemma~\ref{lem:ss-uniform-local-contact},
including $g_1,g_2\ne O$ and $g_2\notin E[3]$. Homogenize its slice
by a fifth parameter $h$ and write
\[
\begin{aligned}
 C&=T+ha,&E_0&=U+hc^2,&D_0&=v+Uc+hd^2+Fa^2,\\
 B&=Tc^2+E_0a,&A&=Td^2+D_0a+Fb,\\
 S&=\begin{pmatrix}A&B&C\\D_0&E_0&h\\F&0&0\end{pmatrix}.
\end{aligned}
\]
These matrices form a projective linear space $\Lambda\simeq\mathbf P^4$
with independent homogeneous coordinates $[F:T:U:v:h]$.
On the affine product put
\[
 X=x+a,\quad Y=y+b,\quad Z=z+c,\quad W=w+d,
 \qquad \ell=w+c^2z+d^2.
\]
Its sections have the equivalent expression
\begin{equation}\label{eq:ss-contact-linear-system}
 s=T\ell+D_0X+E_0Xz+hXw+FY.
\end{equation}
Thus its function basis is $X,Y,\ell,Xz,X\ell$.
The source Hasse cubic restricts to the irreducible cubic
\begin{equation}\label{eq:ss-contact-cubic-homogeneous}
 q=h v^2+h^2v+U^3+h^2U(c^4+c)
                       +F^2T+hF^2(a^4+a)+FT^2.
\end{equation}
Its irreducibility follows by setting $h=1$ and using the odd-pole
Artin--Schreier argument of the preceding lemma; $h$ is not a factor.

\begin{lemma}[A smooth source exists in every nonflex contact slice]
\label{lem:ss-uniform-source-smoothness}
The restriction to $\Lambda$ of the reduced source discriminant has
exactly four forced hyperplane factors, with multiplicities
\[
                            F^3T^3h^3v^2,
\]
and one further irreducible factor, with some positive multiplicity.
The total degree of that remaining factor with multiplicity is $43$.
Consequently the cubic $q=0$ is not contained in the source
discriminant. A geometric generic member of the local family in
Lemma~\ref{lem:ss-uniform-local-contact} is globally smooth.
\end{lemma}
\begin{proof}
We first classify all possible divisorial singularity incidences of
this linear system. Its base points are
\[
              P=(O,O),\qquad Q=(g_1,g_2),
                    \qquad R=(g_1,-2g_2).
\]
The last two are distinct because $g_2$ is not a flex. This follows
on the affine chart from $X=Y=\ell=0$, and on the infinity charts
from the displayed coefficient matrix.

On the affine locus $X\ne0$, ratios of the five basis functions give
\[
                 Y/X,\qquad\ell/X,\qquad z,\qquad\ell.
\]
The second tangent direction is detected by $z$. The first is lost
only when both
\[
                    A_1:=Y+Xx^2=0,\qquad\ell=0.
\]
The first equation has only finitely many points on the first elliptic
curve: it says that the fixed point $g_1$ lies on the tangent at
$(x,y)$, and its displayed nonzero polar polynomial is not identically
zero on the cubic. The second equation has only two geometric points
on the second factor. Thus these exceptions are finite. At each,
first-jet rank is still at least two, since the value and the derivative
of the ratio $z$ are independent. They cannot give a divisorial
singularity incidence in $\mathbf P^4$.

The locus $X=0$ consists of the first-factor fibres through $g_1$
and $-g_1$. On the former the section is $T\ell$. Outside its two
base points, any singular section therefore lies in $T=0$.
At $Q$ the only first derivative is $v$, giving the hyperplane $v=0$.
At $R$ the second derivative direction is $T\partial\ell(R)$,
which is nonzero when $T\ne0$; the two first derivatives at this
base point are independent linear conditions, so its singularity
incidence has codimension two. On the fibre through $-g_1$ one has
$Y=1$, and the value together with the two derivatives have rank
three except when $z=c$, a finite set of points. Their ranks there
are at least two.

On the first-factor fibre through $O$, the section restricts to
$F$ times the section $1$ of $\mathcal O_E(3O)$. Consequently it
is avoided for $F\ne0$ except at $P$; there the only first derivative
is $h$. On the second-factor fibre through $O$, at a finite first
point the value is $T+hX$ and the first derivative is $h$.
The full first-jet rank is three when $X\ne0$; when $X=0$ there
are only two finite exceptional points and rank is at least two.
These calculations cover all four product charts. They show that,
apart from $F,T,h,v$, every divisorial singularity must arise from
the main rank-three jet incidence. That incidence is a
$\mathbf P^1$-bundle over an irreducible open of the surface, and
is therefore irreducible of dimension three.

For clarity its image really is a divisor, and is distinct from all
four hyperplanes. At a moving point satisfying
$X A_1\ell(z+c)\ne0$, solve the three singularity equations as
\[
 F=\frac{T\ell}{A_1},\qquad
 E_0=\frac{T(z^2+c^2)}X+hz^2,\qquad
 D_0=E_0z+hw+Fx^2.
\]
The two remaining parameters are $T,h$. The mixed Hessian is
\[
       \partial_1\partial_2s=E_0+hz^2
                       =\frac{T(z^2+c^2)}X\ne0
\]
when $T\ne0$. Thus the moving point is an ordinary node. The
implicit function theorem for its two first derivatives makes this
node persist only along one local discriminant equation, so the
incidence has divisorial image. Its generic coefficients have
$F,T,h\ne0$. Moreover
\[
 v=\frac{T(z+c)^3}{X}
           +h\bigl(z^2(z+c)+w+d\bigr)+FX^2.
\]
The coefficient of $h$ is a nonzero polar function of the second
elliptic point, so it is generically nonzero. Hence $v$ is also
nonzero at a generic point of this incidence. Its irreducible image,
call it $\Delta_{\mathrm{mov}}$, is distinct from $F,T,h,v$.
The same incidence dimension bounds prove that a generic section
of $\Lambda$ is smooth. In particular restriction of the full
source discriminant is a nonzero homogeneous polynomial.

We compute its four forced multiplicities. At a generic point of
$F=0$, the first-factor fibre through $O$ is a component. Its
residual curve has numerical type $(2,3)$ and meets it at the three
zeros of
\[
                         D_0+E_0z+hw.
\]
These coefficients range freely, so the three zeros are simple,
finite, and avoid the two zeros of $\ell$. The residual curve is
smooth: away from these zeros its first invariant derivative is
nonzero; at a zero the derivative in the second factor is nonzero.
There are therefore exactly three transverse component nodes.
Varying $F$ smooths each to first order, since its restriction at
those finite points is a nonzero local unit. The discriminant order
along $F=0$ is three.

At a generic point of $T=0$, the first-factor fibre through $g_1$
is a component. On its complement the residual equation is
\[
       D_0+E_0z+hw+F\rho=0,\qquad
                     \rho=\frac{y+b}{x+a}.
\]
The rational map $\rho:E\to\mathbf P^1$ has degree two and is
separable, as its derivative $A_1/X^2$ is not identically zero.
It has only finitely many ramification points. For generic
$D_0,E_0,h,F$, no critical value on the second factor coincides
with the image of such a ramification point: after choosing the
other coefficients, only finitely many values of $D_0$ are forbidden.
At the poles the two simple poles of $\rho$ give smooth local
branches. Thus the residual curve is smooth. Its intersections
with the component are the three simple zeros of
\[
                         D_0+E_0z+hw+Fa^2.
\]
They avoid $Q,R$ for a generic choice. Variation of $T$ evaluates
to $\ell$, which is nonzero at all three intersection points.
Hence these nodes are each smoothed to first order and the
multiplicity along $T=0$ is also three.

At a generic point of $h=0$, the only singular point is $P$:
all other possible divisorial incidences have already been shown
to be the distinct hypersurfaces $F,T,v,\Delta_{\mathrm{mov}}$.
In the two $Y=1$ elliptic parameters at $P$ the local equation has
leading terms
\[
          h x+E_0xy+Fy^3+Cx^3+\text{higher terms},
\]
with $E_0,F\ne0$. The Hessian has nonzero mixed coefficient
$E_0$, so the critical point of this equation is unique as a formal
function of $h$. Its critical value has leading term
$Fh^3/E_0^3$: the second coordinate is $h/E_0+O(h^2)$ and the
first is of order at least two. Thus the discriminant has order
exactly three along $h=0$.

At a generic point of $v=0$, its only singularity is similarly the
fixed point $Q$. Choose translated invariant parameters $x,z$ there.
Its quadratic part and transverse parameter have the form
\[
            vx+\alpha x^2+Uxz+T(c^4+c)z^2,
\]
where $U,T(c^4+c)\ne0$. Its critical value therefore has leading
term $T(c^4+c)v^2/U^2$. The discriminant order along $v=0$ is
exactly two. The higher coefficients and the dependence of
$E_0,D_0$ on these transverse parameters do not change the displayed
leading terms, since the critical point has the indicated orders.

The reduced discriminant of the full Segre cubic system has degree
$c_2(J^1L)=3L^2=54$, as proved in the source-discriminant calculation
of Proposition~\ref{prop:ss-level3-three-witness-reduction}.
Its nonzero restriction to $\Lambda$ thus has total degree $54$.
We have exhausted its divisorial support, and found the four orders
$3,3,3,2$. Consequently its remaining irreducible factor
$\Delta_{\mathrm{mov}}$, with its positive multiplicity, has total
degree $54-3-3-3-2=43$.
If the irreducible cubic $q$ were contained in the source discriminant,
it would have to equal $\Delta_{\mathrm{mov}}$, since it is not a
hyperplane. That would make $43$ divisible by three, a contradiction.
Hence a dense open of $q$ parametrizes globally smooth curves.
Intersect it with the nonzero-parameter, simple-contact and unequal
Cartier-value opens of the preceding lemma. Irreducibility of $q$
makes this intersection nonempty, proving the final assertion.
\end{proof}

\begin{remark}[The norm-collision consequences]
This result supplies globally smooth source curves for arbitrary
translations with both coordinates nonzero (interchanging factors
when necessary, and using the checked three-torsion witness in the
finite flex case). Proposition~\ref{prop:ss-uniform-graph-mixed-contact}
now excludes every further nontransverse or triple identification
for odd kernels with injective coordinate projections. The uniform
tangent-triple bridges of Proposition~\ref{prop:ss-uniform-tangent-triple}
then complete the odd-degree Hasse-component argument in
Theorem~\ref{thm:ss-all-odd-Hasse}.
\end{remark}
